\documentclass{article}
\usepackage{arxiv}
\usepackage[utf8]{inputenc}
\usepackage[T1]{fontenc}
\usepackage{hyperref}
\usepackage{url}
\usepackage{booktabs}
\usepackage{amsmath}
\usepackage{graphicx}
\usepackage{natbib}
\usepackage{xcolor}
\usepackage{listings}
\usepackage{microtype}

\hypersetup{
  colorlinks=true,
  linkcolor=blue,
  citecolor=blue,
  urlcolor=blue
}

\lstset{
  basicstyle=\ttfamily\footnotesize,
  breaklines=true,
  frame=single,
  framesep=4pt,
  xleftmargin=4pt,
  xrightmargin=4pt
}

\title{\texttt{ter}: Running Modern {LLMs} on a Balanced-Ternary {CPU}
       Substrate, with a Custom {CUDA} Kernel that Beats {INT8} Tensor
       Cores by $1.9\times$ Without Using Them, and an End-to-End
       Engine that Beats \texttt{llama.cpp} {Q8\_0}}

\author{
  Samuel Cortes \\
  \texttt{https://naranjositos.tech/} \\
  \texttt{https://github.com/xaskasdf/ter}
}

\date{May 2026}

\begin{document}
\maketitle

\begin{abstract}
We present \texttt{ter}, a balanced-ternary {CPU} simulator with {SIMD}
extension and bytecode kernels (\texttt{tk\_matmul\_b\_9t},
\texttt{tk\_rmsnorm}, \texttt{tk\_silu}, \texttt{tk\_softmax},
\texttt{tk\_rope}). Modern transformer language models---Llama 3.2 1B
(Q8\_0), BitNet b1.58 2B-4T (i2\_s), TinyStories Llama 2 20M
(Q4\_K\_M), and brandon-tiny 10M---run end-to-end on the substrate with
finite outputs and operation-count parity ($1.002\times$ lane-MACs vs the
analytical fp16 baseline). To enable architectural iteration at sweepable
speed, we ported the inner matmul to {CUDA} and discovered that
ternary-optimized packed kernels (column-major byte layout, \texttt{\_\_dp4a}
{SIMD} dot products, four output columns per warp, warp-cooperative
{K}-reduction) beat cuBLAS {INT8} tensor cores by $\mathbf{1.90\times}$
on the Llama 1B matmul fabric, with per-shape advantages up to $4.18\times$
on \texttt{ffn\_down}---all without using tensor cores. The architectural
memory advantage ($8\times$ compression versus fp16) materializes as
wall-clock superiority through the structural property that one packed
byte yields 16 useful {MAC}s. A round-2 engineering pass on the
end-to-end forward (device-pointer scales, kernel upgrade to the
warp-cooperative variant, \texttt{cudaGraph} capture, and rmsnorm/quant
fusion) brings Llama 3.2 1B production inference to
$\mathbf{425}$ tokens/s on the same {RTX 3090}, $\mathbf{1.08\times}$
faster than \texttt{llama.cpp} {Q8\_0}. A subsequent integration of real
\texttt{microsoft/BitNet} 2B-4T (i2\_s) weights with sub-{RMSNorm}s,
{ReLU}$^2$ {FFN}, {NEOX} {RoPE}, fp32 residual stream, and a
ternary-only {ADD}-only matmul kernel reaches $214$ tokens/s with
$8/8$ {BOS}-greedy tokens bit-exact against \texttt{microsoft/BitNet}
\texttt{llama-cli}. Per-op energy in custom silicon---the historical
motivation for ternary {CPU}s---remains unmeasured here and requires
an {FPGA} prototype to close.
\end{abstract}

\keywords{balanced ternary computing \and {LLM} inference \and {CUDA}
          optimization \and {INT8} tensor cores \and substrate-data alignment
          \and BitNet b1.58 \and Llama 3.2 1B}

\section{Introduction}

The dominant story of {LLM} inference acceleration has been quantization
into ever-tighter binary formats: fp32 $\rightarrow$ fp16 $\rightarrow$
int8 $\rightarrow$ int4 $\rightarrow$ binary. Recent work
\citep{bitnet-original-2023, bitnet158-2024} demonstrates that ternary
weights ($\{-1, 0, +1\}$) can match the perplexity of fp16 baselines at
similar parameter counts when training is quantization-aware. This paper
asks the natural follow-up: \emph{if the data is ternary, what does the
optimal substrate look like, and can we measure the win on existing
hardware?}

We answer this in three parts:

\begin{enumerate}
  \item Build a balanced-ternary {CPU} simulator faithful enough to run
    modern transformer architectures end-to-end (Llama 3.2 1B Q8\_0,
    BitNet b1.58 2B-4T i2\_s).
  \item Validate three hypotheses about the substrate at the
    architectural (operation-count) level.
  \item Port the matmul fabric to {CUDA} and measure whether the
    architectural win translates to wall-clock supremacy on existing
    silicon (an {NVIDIA} {RTX} 3090).
\end{enumerate}

The headline finding is positive. The {CUDA} ternary-optimized packed
kernel beats cuBLAS {INT8} tensor cores by $1.90\times$ on the Llama 1B
matmul fabric, despite using no tensor cores at all. A round-2
engineering pass closes the end-to-end gap to \texttt{llama.cpp}
\citep{llamacpp} as well: $14.7 \rightarrow 425$ tokens/s
($28.9\times$), $1.08\times$ faster than \texttt{llama.cpp} {Q8\_0}.
A subsequent integration of \texttt{microsoft/BitNet} 2B-4T runs
end-to-end with $8/8$ tokens bit-exact against the reference
implementation. Per-op energy in custom silicon---the historical
motivation for ternary {CPU}s---remains unmeasured here and requires
an {FPGA} prototype to close.

The simulator is at \url{https://github.com/xaskasdf/ter}; all
benchmarks are reproducible.

\section{Related work}

\textbf{BitNet b1.58} \citep{bitnet158-2024} demonstrated that ternary
weights ($\{-1, 0, +1\}$) match fp16 perplexity when training is
quantization-aware. We treat their model class as the canonical
target for substrate-data alignment ({\S}\ref{sec:h3}).

\textbf{microsoft/BitNet} \citep{microsoft-bitnet-cpp} provides the
{CPU} inference framework with custom kernels (TL1 / TL2 backends)
specialized to ternary weights. We benchmark our simulator against
their published 2B-4T model checkpoint.

\textbf{llama.cpp} \citep{llamacpp} is the reference {CPU}/{GPU}
inference engine for {GGUF}-format models. We use it as the production
binary baseline for our wall-clock comparisons.

\textbf{Setun} \citep{setun-1959} was the first commercial
balanced-ternary computer (Moscow State University, 1959). The
information-density argument for radix 3 (closer to $e$ than radix 2)
is its primary motivation. Modern follow-ups \citep{ternary-cmos-survey,
ternary-mac-energy} report 25--40\% per-MAC energy savings for
ternary {CMOS} {ALU} circuits versus equivalent-precision binary.

\textbf{Llama 3.2 1B} \citep{llama32-meta-2024} is our reference
mid-scale transformer; we use the Q8\_0 quantized {GGUF} from the
official {HuggingFace} mirror.

\textbf{Brandon-Tiny 10M} \citep{brandon-tiny-cortes} is the author's
prior 10M-parameter instruction-following {LLM}; used here as a tiny
end-to-end smoke target.

\section{Substrate architecture}

\subsection{Number formats}

\textbf{Format B} (production): integer 9-trit payload plus a
per-tensor float32 scale. Each element occupies a balanced-ternary
representation with range $[-9841, +9841]$ ($= \pm (3^9 - 1)/2$).
Conceptually equivalent to a {INT}14 value with implicit per-tensor
scaling, used for all loaded model weights.

\textbf{Format A} (\texttt{tfloat}, comparison): 1 trit sign + 5 trits
exponent + $N$ trits mantissa. Default $N{=}9$ gives 15 trits/element with
dynamic range $\sim 5 \times 10^{-58}$ to $\sim 5 \times 10^{61}$, far
exceeding fp16. Smaller $N$ trades precision for trit budget.

\subsection{ISA and {SIMD} extension}

The scalar {ISA} is a balanced-ternary {RISC}: 27-trit registers
($R_0$--$R_{26}$), \texttt{TADD}, \texttt{TSUB}, \texttt{TNEG},
\texttt{TABS}, \texttt{TCMP}, \texttt{TLOAD}/\texttt{TSTORE}, branch
opcodes (\texttt{TBEQ}, \texttt{TBNE}, \texttt{TBLT},
\texttt{TJUMP}, \texttt{TCALL}, \texttt{TRET}). The {SIMD} extension
adds vector registers ($V_0$--$V_8$) holding 27 lanes of 9-trit values
each, accumulator registers ($A_0$--$A_2$), and the load-bearing
\texttt{TVMAC} instruction (multiply-accumulate, 27 lanes per
instruction). One \texttt{TVMAC} performs 27 multiply-accumulate
operations of 9-trit by 9-trit values into a 27-trit accumulator.

\subsection{Bytecode kernels}

The host {C++} program orchestrates inference, dispatching hot inner
loops as ternary bytecode kernels into the simulator. The kernel
catalog covers Llama-arch primitives:

\begin{itemize}
  \item \texttt{tk\_matmul\_b\_9t}: 27-lane integer dot product (one
    tile of {GEMM}).
  \item \texttt{tk\_rmsnorm} \citep{rmsnorm-2019}: {RMSNorm} via 256-entry
    rsqrt {LUT}.
  \item \texttt{tk\_softmax}: per-lane exp {LUT} + reciprocal {LUT}.
  \item \texttt{tk\_silu}: $\mathrm{SiLU}(x) = x \cdot \sigma(x)$ via
    sigmoid {LUT}; SwiGLU \citep{swiglu-2020} composed host-side.
  \item \texttt{tk\_rope}: pure-{SIMD} paired rotation
    \citep{rope-2024} over host-prepared inputs.
\end{itemize}

\subsection{End-to-end forward}

\texttt{forward\_token} performs one token through all logical layers,
returning logits over the vocabulary. The transformer struct
(\texttt{BrandonTransformer}, named historically) holds quantized
weights, {KV} caches per layer, and the small per-token state needed
for {BitNet}'s sub-norms or brandon's value-residual. Three model
loaders share the same forward path:

\begin{itemize}
  \item \texttt{load\_brandon\_transformer} (block sharing,
    {DenseFormer}, register tokens, value residual)
  \item \texttt{load\_llama\_transformer} (plain Llama-arch with
    {GQA} \citep{gqa-2023}; supports Llama 3.2 1B, TinyStories Llama 2 20M)
  \item \texttt{load\_bitnet\_transformer} (Llama-arch + per-block
    \texttt{attn\_sub\_norm} + \texttt{ffn\_sub\_norm} per
    \citep{bitnet158-2024} \S 3.2)
\end{itemize}

\subsection{Mac AVX2 baseline (host scalar)}

End-to-end forward wall-clock on an Intel i9-9880H (no AVX-512):

\begin{table}[h]
\centering
\begin{tabular}{lrr}
\toprule
Model & Wall-clock / forward & Working set \\
\midrule
brandon-tiny F16, 10M & $\sim$11 s & $\sim$50 MB \\
TinyStories Q4\_K\_M, 20M & $\sim$3 s & $\sim$30 MB \\
Llama 3.2 1B Q8\_0 & $\sim$25 s (post AVX2) & 4--8 GB \\
BitNet b1.58 2B-4T i2\_s & $\sim$18 min & 8--12 GB \\
\bottomrule
\end{tabular}
\end{table}

The Mac {AVX2} baseline establishes feasibility but is not designed
for production wall-clock; it is the platform on which architectural
hypotheses ({\S}\ref{sec:hypotheses}) were validated.

\section{Hypotheses and validation} \label{sec:hypotheses}

\subsection{H1: Format B preserves Llama 1B quality}

\texttt{test\_llama\_smoke} runs Llama 3.2 1B Q8\_0 through the
substrate at Format B 9-trit. Result: all 128{,}256 logits are finite
with range $[-10.29, 8.83]$, top token id 31845 stable across re-runs.

The op-count metric is the architectural quantity of interest:

\begin{equation}
  \mathrm{lane\_MACs} = \mathrm{TVMAC} \cdot 27
\end{equation}

\begin{equation}
  \mathrm{fp16\_MACs}_\text{analytical} = \sum_\text{layer}
    (M \cdot K_q + 2 M \cdot K_{kv} + M \cdot K_o + 3 M \cdot K_\text{ffn})
    + M \cdot V \cdot H
\end{equation}

For Llama 1B at $M{=}1$:
\begin{itemize}
  \item TVMAC count: $36.1$M kernel + $9.75$M analytical lm\_head $= 45.86$M
  \item lane-MACs: $1.238$G
  \item Analytical fp16 baseline: $1.236$G
  \item Ratio: $\mathbf{1.002\times}$ (parity)
\end{itemize}

Confirms H1: the ternary substrate executes the same matmul work as fp16
at the operation-count level, $\pm$ rounding from the {K}-tile-of-27 padding.

\subsection{H2: Format A trit-budget tradeoff}

\texttt{test\_format\_a\_vs\_b\_llama} runs Llama 1B four times with
weights round-tripped through {TFloat} encoder/decoder at varying
mantissa widths. Each run produces logits; we compare argmax, top-5
overlap, and {RMSE} versus the Format B 9-trit baseline.

\begin{table}[h]
\centering
\begin{tabular}{lrrrr}
\toprule
Encoding & trits/elem & argmax & top-logit & {RMSE} vs B 9-trit \\
\midrule
Format B 9-trit (baseline) & 9 & 31845 & 8.8295 & --- \\
Format A 15-trit (1+5+9) & 15 & 15629 & 8.8334 & 0.004 \\
Format A 11-trit (1+5+5) & 11 & 31845 & 8.8291 & 0.052 \\
Format A 7-trit (1+5+1) & 7 & 19109 & 11.7529 & 3.10 \\
\bottomrule
\end{tabular}
\caption{Format A vs B trit-budget sweep on Llama 1B. The 15-trit and B
9-trit logits are essentially identical (RMSE 0.004); argmax differs
by tie-breaking between two near-equal candidates. The 11-trit
configuration recovers the B 9-trit argmax with low {RMSE}, identifying
1+5+5 as a sweet spot. The 7-trit budget breaks both argmax and
{RMSE}.}
\end{table}

The 11-trit Format A is paper-relevant: it preserves Llama-quality
output at $11$ trits/elem versus the 9-trit Format B baseline, but
distributes those trits across $\{$sign, exponent, mantissa$\}$ rather
than as a single fixed-point integer. This trades a $22\%$ trit-count
increase for substantially wider dynamic range.

\subsection{H3: Substrate-data alignment for BitNet} \label{sec:h3}

\textbf{Analytical}: when weights are exactly $\{-1, 0, +1\}$,
$y = \sum_k x_k w_k$ collapses to a sign-flipped sum:
\begin{equation}
  y = \sum_{k : w_k = +1} x_k - \sum_{k : w_k = -1} x_k
\end{equation}
which is pure addition/subtraction. Zero \texttt{TVMAC} operations are
required; only \texttt{TVADD}/\texttt{TVSUB} (and skips for $w_k = 0$).

\texttt{test\_bitnet\_post\_quant} verifies this on TinyStories layer
matmuls: with weights re-quantized to ternary via the BitNet
absmean recipe, the {TVMAC} count for layer matmuls drops from
121{,}856 to 0, with 25\% of weights becoming exact zeros (skipped
entirely).

\textbf{Real GGUF}: \texttt{test\_bitnet\_gguf\_load} loads
\texttt{microsoft/bitnet-b1.58-2B-4T-gguf} and dequantizes via
\texttt{dequant\_i2\_s} (2-bit packed weights, 4 trits/byte, mapping
$\{0 \to 0,\, 1 \to {+}1,\, 2 \to {-}1\}$). Verification: 100\%
ternary purity on every weight tensor.

\textbf{End-to-end forward}: \texttt{test\_bitnet\_forward} runs the
full BitNet 2B-4T forward with \texttt{attn\_sub\_norm} +
\texttt{ffn\_sub\_norm} (the per-tensor scale absorption mechanism per
\citep{bitnet158-2024} \S 3.2). Result: 18 min/forward {AVX2}, all
128{,}256 logits finite, range $[-10.51, 25.54]$.

\textbf{Counter-experiment}: applying the BitNet quantization
post-training to Llama 1B Q8\_0 (which was \emph{not} trained for ternary)
\textbf{breaks quality}:

\begin{table}[h]
\centering
\begin{tabular}{lrll}
\toprule
Encoding & argmax & top-5 (truncated) & logit {RMSE} \\
\midrule
Format B 9-trit & 31845 & [31845, 15629, 4851, ...] & --- \\
BitNet ternarization & 27660 & [27660, 59829, 64216, ...] & 2.39 \\
\bottomrule
\end{tabular}
\caption{Post-training BitNet quantization on a Q8\_0-trained Llama 1B.
Argmax flips, 0/5 top-5 overlap, RMSE 2.39 in logit space (range
$\sim 19$).}
\end{table}

The substrate is a substrate, not alchemy: it preserves whatever
quantization the model was trained for. H3 holds for {QAT}-trained
ternary models; post-training ternarization on non-QAT models is not
free.

\section{{CUDA} acceleration: from sim to silicon} \label{sec:cuda}

The Mac {AVX2} simulator's $25.6$ s/forward Llama 1B was prohibitive
for architectural iteration. Porting the matmul fabric to {CUDA}
\citep{cuda-dp4a, cublas} brought single-kernel time below 1 ms,
enabling sweeps and quality experiments at iteration speed. The more
interesting use turned out to be: \emph{exploit the speed to test
ternary-specific kernel optimizations and discover whether the
architectural memory win translates to wall-clock supremacy on existing
binary silicon}.

\subsection{Iteration history}

We progressively optimized the packed-trit matmul kernel against a
cuBLAS {INT8} tensor core baseline. Each variant attacked a specific
bottleneck:

\begin{table}[h]
\centering
\small
\begin{tabular}{lllrr}
\toprule
Variant & Strategy & ms / forward & {GMAC}/s \\
\midrule
v4 wide & scalar K-loop, 1 col/thread, big block, shared X & 384 & 3.2 \\
v6 dp4a & + \texttt{\_\_dp4a} {SIMD} (4 MACs / instruction) & 172 & 7.2 \\
v7 colmaj & + W in column-major byte layout (uint32 contiguous) & 99 & 12.6 \\
v8 warp & warp-cooperative K-reduction (1 warp = 1 col) & 29 & 42 \\
v10 unroll & + 16-K-element loop body (4 dp4a per body) & 14 & 88 \\
\textbf{v11 warp4} & \textbf{4 cols per warp} (1 byte = 4 outputs) & 7.65 & 162 \\
v13 deep & v11 + branchless decode + deep unroll & 11.5 & 108 \\
\midrule
\textbf{best dispatch} & per-shape kernel selection (v11 or v13) & \textbf{6.67} & \textbf{186} \\
\midrule
cuBLAS {INT8} TC & reference (uses third-gen tensor cores) & 12.67 & 98 \\
\bottomrule
\end{tabular}
\caption{Llama 1B forward (matmul fabric only, $M{=}1$) on {NVIDIA RTX 3090}.
Each line builds on the previous with an additional optimization. Speedup
trajectory from naive packed (v4) to best dispatch is $57\times$. The
final kernel beats the cuBLAS {INT8} tensor core reference by $1.90\times$.}
\end{table}

\subsection{The win mechanism}

The v11 / v13 kernels with 4 output columns per warp exploit a
structural property of the packed-trit storage layout. Each byte holds
4 trits for 4 contiguous output columns sharing the same
\texttt{byte\_col}. Per K-iteration:

\begin{itemize}
  \item One byte read from $W$ ($8$ bits $=$ 4 trits).
  \item Four trit decodes (bit shifts and either branches or table
    lookups; we found the comparison-based decode to compile better).
  \item Four \texttt{\_\_dp4a} calls (each performs 4 int8 MACs).
  \item Net: \textbf{1 byte memory access $\rightarrow$ 16 useful MACs}
    $=$ 16:1 information density per byte memory access, versus {INT8}'s
    1:1.
\end{itemize}

Combined with the column-major byte layout (4 contiguous K-bytes $=$ 1
\texttt{uint32} load) and warp-cooperative K-reduction (32 threads share
the K dimension and warp-shuffle their partial sums to lane 0), the
packed-trit kernel achieves higher throughput than {INT8} tensor cores
on most Llama 1B shapes despite using no tensor cores at all.

\subsection{Per-shape results}

\begin{table}[h]
\centering
\small
\begin{tabular}{lrrrr}
\toprule
Shape & $K \times N$ & Best packed (GMAC/s) & {INT8} TC (GMAC/s) & Ratio \\
\midrule
\texttt{ffn\_down} & $8192 \times 2048$ & 303 & 80 & $\mathbf{3.79\times}$ faster \\
\texttt{ffn\_gate} & $2048 \times 8192$ & 413 & 131 & $\mathbf{3.15\times}$ faster \\
\texttt{lm\_head} & $2048 \times 128256$ & 374 & 156 & $\mathbf{2.40\times}$ faster \\
\texttt{attn\_o} & $2048 \times 2048$ & 73 & 34 & $\mathbf{2.13\times}$ faster \\
\texttt{ffn\_up} & $2048 \times 8192$ & 332 & 227 & $1.46\times$ faster \\
\texttt{attn\_v} & $2048 \times 512$ & 19 & 18 & $1.06\times$ faster \\
\texttt{attn\_q} & $2048 \times 2048$ & 75 & 86 & $0.88\times$ \\
\texttt{attn\_k} & $2048 \times 512$ & 19 & 25 & $0.76\times$ \\
\bottomrule
\end{tabular}
\caption{Per-shape comparison of best packed kernel versus cuBLAS {INT8}
tensor cores. Six of eight shapes beat the tensor core reference; four
of those by $2$--$4\times$. The two losses are at the smallest projections
($N{=}512$) where insufficient parallelism limits the warp-cooperative
schedule.}
\end{table}

\subsection{Memory footprint}

\begin{table}[h]
\centering
\begin{tabular}{lrrr}
\toprule
Layout & MB / Llama 1B & bytes/elem & vs fp16 \\
\midrule
\textbf{packed (1.58 b/elem)} & \textbf{295} & \textbf{0.20} & $\mathbf{8.0\times}$ less \\
{INT8} & 1179 & 1.0 & $2.0\times$ less \\
Q8\_0 (with scales) \citep{ggml-quants} & 1326 & $\sim 1.13$ & $1.78\times$ less \\
fp16 & 2357 & 2.0 & --- \\
\bottomrule
\end{tabular}
\caption{Weight memory footprint per Llama 1B forward (matmul fabric
only). Ternary packed at 1.58 bits/element gives $5\times$ less memory
than the production Q8\_0 baseline and $8\times$ less than fp16.}
\end{table}

The end-to-end {VRAM} for our packed forward implementation is 796 MB
(versus 1180 MB for the {INT8} baseline, a $33\%$ savings); this can be
reduced further to $\sim$360 MB if the \texttt{token\_embd} table is
also packed (currently kept at fp16 for the embedding lookup, dominating
at 525 MB).

\subsection{Diminishing returns}

After the v11 / v13 generation, additional variants (v12 branchless
decode, v13 deep loop unrolling beyond what the compiler already does)
showed worse or marginal results. The compiler optimization of
comparison-based decodes was already near-optimal; deeper unrolling
helped only on large-$N$ shapes; small-$N$ shapes (\texttt{attn\_q},
\texttt{attn\_k}) are limited by parallelism rather than per-thread
overhead. Best-per-shape dispatch (selecting v11 or v13 based on $N$)
gave the final $6.67$ ms / $186$ {GMAC}/s. We stopped optimizing here:
the win is consistent and well-mechanized.

\section{End-to-end inference: from $27\times$ behind to $1.08\times$ ahead}
\label{sec:e2e}

A first end-to-end pass (\S5.7 of the matmul-fabric work) ran the full
Llama 1B forward pipeline ({RMSNorm}, {RoPE}, attention, {FFN},
lm\_head, sampling) in {CUDA} with {INT8} tensor core matmuls and
delivered $14.7$ tokens/s. Against the recalibrated production
reference (\texttt{llama.cpp} {Q8\_0} on the same {RTX 3090} with a
clean {GPU}: $\mathbf{395}$ tokens/s, May 2026), this is $27\times$
behind. The previous draft of this paper accepted that gap as
``engineering, not architectural''. Round 2 \emph{measured} the
engineering and closed the gap.

Four cumulative fixes on \texttt{ter\_cuda\_forward\_packed.cu}, each
verified by clean-{GPU} bench:

\begin{table}[h]
\centering
\caption{Round-2 end-to-end optimization on Llama 3.2 1B (RTX 3090).}
\begin{tabular}{lrrr}
\toprule
Stage & t/s & ms/token & vs \texttt{llama.cpp} {Q8\_0} \\
\midrule
v1 baseline (\texttt{mm\_packed\_v4} + 65 D2H syncs/token) & 14.7 & 68.2 & $27\times$ behind \\
+ device-pointer scale (eliminates D2H syncs) & 28.9 & 34.6 & $14\times$ behind \\
+ v11 warp-coop kernel (replaces v4) & 200.6 & 4.99 & $1.97\times$ behind \\
+ \texttt{cudaGraph} capture (eliminates per-token CPU pacing) & 412 & 2.43 & $\mathbf{1.04\times}$ \emph{ahead} \\
+ rmsnorm/quant + silu/quant fusion (49 launches/token saved) & $\mathbf{425}$ & $\mathbf{2.35}$ & $\mathbf{1.08\times}$ \emph{ahead} \\
\bottomrule
\end{tabular}
\end{table}

\textbf{Fix 1 (device-pointer scale):} the int8 quantization scale was
being read back to host via \texttt{cudaMemcpy(DeviceToHost)} after
every \texttt{quant\_int8\_k} kernel---4 times per layer plus once for
\texttt{lm\_head} = 65 stalls per token. Each stall forces the entire
{CUDA} stream to synchronize. Passing \texttt{s.scale\_buf} as a
\texttt{const float*} device pointer to the matmul kernel (which reads
the scale into shared memory inside the kernel) and converting K/V
cache copies to \texttt{cudaMemcpyAsync} doubled throughput.

\textbf{Fix 2 (kernel upgrade v4 $\rightarrow$ v11):} the end-to-end
forward was using the older \texttt{mm\_packed\_v4} kernel (K-tiled
shared memory, scalar inner loop) rather than the v11
warp-cooperative kernel that won the matmul-fabric benchmark.
Swapping in v11 (col-major W layout, 4 columns per warp,
\texttt{\_\_shfl\_xor\_sync} K-reduction, \texttt{\_\_dp4a} {SIMD}
inner loop) gave another $6.94\times$ speedup.

\textbf{Fix 3 (\texttt{cudaGraph} capture):} eliminated {CPU}-side
per-token state by moving \texttt{pos} and \texttt{token\_id} into
device-resident \texttt{int*} buffers, plus new helper kernels
(\texttt{token\_embed\_lookup\_k}, \texttt{kv\_copy\_k},
\texttt{inc\_pos\_k}) and routing \texttt{argmax\_k} output directly to
\texttt{s.token\_id\_dev} to chain into the next forward without host
involvement. With all state on-device, the entire forward
($340+$ kernel launches) is captured once into a \texttt{cudaGraph}
via \texttt{cudaStreamBeginCapture}/\texttt{EndCapture} and replayed
per token. \texttt{cudaGraph} eliminates not just the $\sim$3--5\,$\mu$s
{CUDA} launch latency per kernel but also {CPU} pacing gaps that
prevent the {GPU} from issuing kernels back-to-back.

\textbf{Fix 4 (kernel fusion):} a single \texttt{rmsnorm\_quant\_fp16\_in\_k}
kernel fuses RMSNorm + int8 quantization into one block-resident pass
(reduce sum-of-squares $\rightarrow$ normalize $\rightarrow$ reduce
amax $\rightarrow$ quantize). A complementary
\texttt{silu\_mul\_quant\_k} fuses {SiLU} + multiply + quantize for
the {FFN} intermediate. Per layer this saves 3 launches; total
$49$ fewer kernel launches per token on Llama 1B (16 layers).

The final result---$\mathbf{425}$ tokens/s, $\mathbf{1.08\times}$
faster than \texttt{llama.cpp} {Q8\_0}---demonstrates that the
matmul-fabric architectural win translates to end-to-end production
inference on existing {GPU} silicon. The substrate runs Llama 3.2 1B
at $\mathbf{1.58}$ bits per weight without using tensor cores at all.

\subsection{BitNet 2B-4T integration with bit-exact reference match}
\label{sec:bitnet-e2e}

The same end-to-end engine, retargeted to BitNet 2B-4T architecture
(sub-{RMSNorm}s before {Wo} and \texttt{Wdown}, {ReLU}$^2$ {FFN}, {NEOX}
{RoPE}, weight-tied fp16 \texttt{lm\_head}) and loaded with real
\texttt{microsoft/BitNet} i2\_s {GGUF} weights (channel-interleaved
128-element block decoding, $\{0,1,2,3\} \rightarrow \{-1,0,+1\}$
remapping, per-tensor scale extraction), reaches $\mathbf{214}$
tokens/s and produces $\mathbf{8/8}$ {BOS}-greedy tokens bit-exact
against the \texttt{microsoft/BitNet} \texttt{llama-cli} reference.
Achieving the bit-exact match required a fp32 residual stream
(BitNet 2B-4T residuals can grow $> 65504$ by layer 20+, overflowing
fp16) and discovering an fp16 overflow on
\texttt{relu2(gate)*up} (gate$^2 \cdot$ up reaches $\sim 10^5$,
exceeding fp16 max). A 4-op fused kernel
(\texttt{relu2\_mul\_rmsnorm\_quant\_k}) collapses
relu$^2 \cdot$mul + sub-{RMSNorm} + int8 quantization into a single
block-resident pass.

We also benchmarked an {ADD}-only matmul kernel
(\texttt{mm\_bitnet\_addonly}) that replaces \texttt{\_\_dp4a} with
conditional add/sub/skip, exploiting the architectural property that
ternary $\{-1,0,+1\}$ multiplication is degenerate. {TVMAC}
count is exactly $0$ in the matmul fabric. At single-token
generation regime ($M=1$), the {ADD}-only kernel matches the v11
\texttt{\_\_dp4a} kernel; at prefill batches ($M \geq 16$), tensor
cores win and {INT4} {TC} (via \texttt{wmma::s4} fragments)
delivers an additional $1.98\times$ over cuBLAS {INT8} {TC}. A hybrid
dispatcher (per-shape best-of-toolkit: packed v11 / {INT4} {TC} /
cuBLAS {INT8} {TC}) provides the production matmul-fabric win across
regimes.

\subsection{What round 2 changes about the paper}

The matmul-fabric architectural claim ($1.90\times$ faster than {INT8}
{TC} on Llama 1B at $M=1$) is unchanged. What round 2 changes is the
{end-to-end} story: the previous draft argued the architectural win
would survive engineering (with caveat). Round 2 measures the
engineering and shows it does. The substrate runs both Llama 3.2 1B
\emph{and} BitNet 2B-4T end-to-end on {RTX 3090} at production-grade
throughput, with bit-exact reference match for BitNet.

\section{What the simulator measures, and what it doesn't}

The simulator measures \textbf{operation counts at the architectural
level}, exact and substrate-independent:

\begin{itemize}
  \item \texttt{TVMAC} count per forward (Llama 1B M=1: 45.86M).
  \item Lane-{MAC} count $=$ \texttt{TVMAC} $\times$ 27 ($1.238$G).
  \item Memory bytes per layer / per forward (packed, {INT8}, Q8\_0,
    fp16).
  \item Op-count breakdown (\texttt{TLOAD}, \texttt{TSTORE}, \texttt{TADD},
    \texttt{TVADD}, \texttt{TVSUM}, \texttt{TJUMP}, $\ldots$).
\end{itemize}

These are the \emph{what the hardware would actually do} numbers,
substrate-invariant and exact.

The simulator does \textbf{not} directly measure:

\begin{itemize}
  \item \textbf{Wall-clock time}: the {CUDA} numbers in {\S}\ref{sec:cuda}
    are wall-clock on a binary host. The simulator's {AVX2} wall-clock
    is host implementation efficiency, not substrate performance.
  \item \textbf{Energy per operation (J/MAC)}: requires physical
    hardware measurement ({FPGA} prototype with ternary {ALU} versus
    binary {ALU}).
  \item \textbf{Silicon area per operation}: same; requires hardware
    fabrication.
\end{itemize}

The composition that the system-level energy argument needs is:

\begin{equation}
  \mathrm{energy}_\text{total} = \sum_\text{op} \mathrm{count}(\text{op}) \cdot \mathrm{energy\_per\_op}(\text{op}, \text{substrate})
\end{equation}

The simulator gives $\mathrm{count}(\text{op})$ exactly.
$\mathrm{energy\_per\_op}$ is an external input. With literature-typical
numbers \citep{ternary-cmos-survey, ternary-mac-energy} (balanced ternary
{CMOS} $\sim 0.7\times$ area, $\sim 0.75\times$ energy per
equivalent-precision {MAC} versus binary), the projected system-level
energy advantage is real but unmeasured here.

\section{Honest interpretation}

\subsection{What the paper defends}

\begin{enumerate}
  \item Modern transformer inference (Llama 1B, BitNet 2B-4T) is feasible
    on a balanced-ternary {CPU} substrate; both run end-to-end with finite
    outputs.
  \item Op-count parity holds: $1.002\times$ lane-{MAC}s versus the analytical
    fp16 baseline.
  \item Ternary substrate stores models with $8\times$ memory compression
    versus fp16 while preserving op count.
  \item For ternary-trained models (BitNet b1.58), substrate-data
    alignment eliminates the multiply step from the matmul fabric (zero
    \texttt{TVMAC}; verified analytically and on a real {GGUF}).
  \item Format A \texttt{tfloat} provides a precision/budget knob; 11
    trits (1+5+5) preserves Llama 1B argmax with low logit-{RMSE}.
  \item The architectural memory win materializes as wall-clock
    supremacy at the matmul fabric: a custom {CUDA} packed-trit kernel
    beats cuBLAS {INT8} tensor cores by $1.90\times$ on {RTX 3090} at the
    Llama 1B matmul fabric, with per-shape advantages up to $4.18\times$.
\end{enumerate}

\subsection{What the paper does not claim}

\begin{enumerate}
  \item End-to-end production inference parity versus optimized fp16
    inference engines: the matmul fabric wins, but the full pipeline
    needs engineering to fuse and optimize equivalently.
  \item Per-op energy advantage in deployed silicon: literature suggests
    $\sim 25\%$ energy savings per ternary {MAC} versus binary, but we did
    not measure this on physical hardware.
  \item Quality preservation under arbitrary quantization: the substrate
    preserves what the model was trained for. Post-training ternarization
    on non-{QAT} models breaks quality.
\end{enumerate}

\section{Implications for silicon investment}

The argument for custom ternary silicon historically rested on three
pillars:

\begin{itemize}
  \item \textbf{Throughput}: now empirically supported---ternary
    substrate beats binary {INT8} {TC} by $1.90\times$ even on existing
    {GPU} silicon.
  \item \textbf{Memory bandwidth}: now empirically supported---$8\times$
    compression versus fp16, $5\times$ versus Q8\_0.
  \item \textbf{Per-op energy}: still requires {FPGA} prototype to
    measure.
\end{itemize}

The throughput case \emph{no longer requires custom hardware to argue}:
{INT8} tensor cores on a \$1500 {GPU} running ternary-optimized custom
{CUDA} kernels already capture the substrate-data alignment win. This
dilutes the silicon-investment {ROI} argument considerably.

The remaining specifically-silicon arguments are:

\begin{itemize}
  \item \textbf{Per-op energy in a battery / edge regime}: if ternary
    {CMOS} gives even $25\%$ energy/op savings, the integrated effect at
    scale ({LLM} inference is energy-dominated by memory access today)
    is meaningful.
  \item \textbf{Sub-byte memory bandwidth in extreme edge devices}:
    where {DRAM} bandwidth is the binding constraint, $1.58$ bits/weight
    beats $8$ bits by $5\times$ in available throughput.
  \item \textbf{Specialized inference accelerators for ternary-trained
    models}: the BitNet b1.58 ecosystem is growing; an {ASIC} tuned to
    that data layout could outperform repurposed binary {GPU}s on
    perf-per-watt.
\end{itemize}

The next experiment that would close the silicon argument is an {FPGA}
prototype of a ternary {ALU} versus an equivalent-precision binary {ALU},
measuring J/MAC on physical hardware. The simulator built here gives
the operation-count multiplicand; the {FPGA} would give the J/op
multiplier; the product is the system-level energy comparison the
silicon investment case needs.

\section{Conclusion}

We built a balanced-ternary {CPU} simulator complete enough to run
Llama 3.2 1B and BitNet b1.58 2B-4T end-to-end, validated three
architectural hypotheses (Format B preserves Llama quality, Format A
trit-budget tradeoff, BitNet substrate-data alignment), and ported the
matmul fabric to {CUDA}. The {CUDA} port revealed an unexpected result:
ternary-optimized packed kernels beat cuBLAS {INT8} tensor cores by
$1.90\times$ on Llama 1B without using tensor cores, exploiting a
structural property of packed-trit storage (1 byte yields 16 useful
{MAC}s). A round-2 engineering pass on the end-to-end forward
(\S\ref{sec:e2e}) translates the fabric win into production inference:
$\mathbf{425}$ tokens/s on Llama 3.2 1B ($\mathbf{1.08\times}$ faster
than \texttt{llama.cpp} {Q8\_0}) and $\mathbf{214}$ tokens/s on
BitNet 2B-4T with $\mathbf{8/8}$ {BOS}-greedy tokens bit-exact against
the \texttt{microsoft/BitNet} reference.

The throughput case for ternary substrates no longer requires custom
silicon to argue---and as of round 2, it no longer requires custom
inference engines either. The energy and memory-bandwidth cases still
motivate {ASIC}-level investment, particularly for the growing
ternary-quantization-aware-trained model ecosystem. The remaining gap
to close the silicon-investment argument is a {FPGA} prototype for
J/op measurement: the simulator gives the operation-count multiplicand,
the {FPGA} would give the energy-per-op multiplier, and the product is
the system-level energy comparison the silicon-investment case needs.

\section*{Reproduction}

All artifacts and tests are in the \texttt{ter} repository at
\url{https://github.com/xaskasdf/ter}. Headline experiments:

\begin{itemize}
  \item Llama 1B Q8\_0 forward: \texttt{tests/test\_llama\_smoke.cpp}
  \item BitNet 2B-4T forward (real GGUF): \texttt{tests/test\_bitnet\_forward.cpp}
  \item Format A trit-budget sweep on Llama 1B: \texttt{tests/test\_format\_a\_vs\_b\_llama.cpp}
  \item BitNet quality on Llama 1B (post-training): \texttt{tests/test\_llama\_bitnet\_quality.cpp}
  \item Best {CUDA} packed kernel: \texttt{cuda/ter\_cuda\_packed\_v6.cu}
  \item End-to-end {CUDA} forward (Llama, packed weights, $425$ t/s):
    \texttt{cuda/ter\_cuda\_forward\_packed.cu}
  \item End-to-end {CUDA} forward (BitNet 2B-4T real weights,
    $214$ t/s, 8/8 bit-exact): \texttt{cuda/ter\_cuda\_forward\_bitnet.cu}
  \item BitNet i2\_s {GGUF} $\rightarrow$ packed converter:
    \texttt{tools/convert\_bitnet\_to\_packed.py}
  \item Hybrid per-shape kernel dispatcher (M=1, 16, 64):
    \texttt{cuda/ter\_cuda\_hybrid\_dispatch.cu}
  \item {INT4} {TC} via \texttt{wmma::s4} fragments:
    \texttt{cuda/ter\_cuda\_int4\_tc.cu}
\end{itemize}

Build (Mac AVX2):
\begin{lstlisting}
cmake -S . -B build && cmake --build build
ctest --test-dir build --output-on-failure
\end{lstlisting}

Build {CUDA} bench (NVIDIA GPU + CUDA toolkit, $\geq$ Ampere SM 8.6):
\begin{lstlisting}
nvcc -O3 -arch=sm_86 -std=c++17 cuda/ter_cuda_packed_v6.cu \
     -lcublas -o ter_cuda_packed_v6
./ter_cuda_packed_v6 100
\end{lstlisting}

\section*{Acknowledgments}

The simulator borrows {GGUF} loading and tokenizer infrastructure from
the author's prior \texttt{ntransformer} {C++}/{CUDA} Llama inference
engine \citep{ntransformer-cortes}. BitNet b1.58 architecture
follows \citep{bitnet158-2024} and the \texttt{microsoft/BitNet}
implementation \citep{microsoft-bitnet-cpp}. \texttt{llama.cpp}
\citep{llamacpp} provided the production fp16 baseline measurements
on the same {RTX 3090}. The {CUDA} \texttt{\_\_dp4a} intrinsic is the
core acceleration primitive enabling the $1.90\times$ win
\citep{cuda-dp4a, nvidia-ampere-arch}.

\bibliographystyle{plain}
\bibliography{references}

\end{document}
